\section{Experimental Methodology}

\subsection{Objectives of the Evaluation}
The evaluation is designed as a system-level validation campaign rather than as a micro-benchmark study. The goal is to verify whether the complete RETROSPECT stack behaves coherently under realistic deployment and supervision conditions while also producing a quantitative profile suitable for a Transactions-style systems paper. In particular, the evaluation seeks to answer four questions: can a device securely enroll through the gateway; can liveness remain observable and stable during periodic heartbeat traffic; can application deployment complete with positive runtime acknowledgment on the target device; and can the whole workflow be reproduced with statistically stable timing and success-rate evidence on a commodity research platform.

This framing is deliberate. Because the paper's contribution is the integration of orchestration, protocol handling, and embedded execution, the most relevant evidence is end-to-end functional correctness under operational conditions. For the same reason, the study reports not only positive-path latency, but also transactional success, percentile behavior, and data-source provenance for each metric so that reviewers can distinguish actual runtime success from stale or noisy control-plane state.

\subsection{Evaluation Environment}
The experiments are executed on a single-node K3s cluster hosting the RETROSPECT services in the dedicated namespace \texttt{wasmbed}. The deployed stack includes the API server, gateway, CRDs, and the auxiliary resources needed to drive device enrollment and status management. A local image registry is used to build and distribute container images inside the test environment, enabling fast iteration while preserving realistic cluster deployment semantics. During the quantitative campaign, the application controller is deployed with one active replica. Its reconciliation loop has been hardened to detect and repair stale CRD status, making the Application CRD \texttt{phase} field the authoritative signal for deployment success. This allows the measurements to cover the full control-plane path, from API intent through gateway-mediated device execution to northbound CRD phase reconciliation.

On the southbound side, the evaluation uses a Renode-emulated Zephyr device with a deterministic identity and a known public key. The firmware establishes a TLS session with the gateway, exchanges CBOR-framed enrollment messages, and then emits periodic heartbeats. The experiment harness reaches the cluster through controlled port-forwarding and TAP-based network setup, which keeps the setup faithful to the intended embedded target profile while still supporting reproducible regression-oriented validation.

\subsection{Experimental Matrix and Measured Variables}
The validation campaign is divided into five experiment classes. Each class is measured independently, and each trial produces a raw record together with a machine-readable summary. The default campaign consists of one warm-up run followed by 50 independent trials per condition. With $n=50$, the Wilson score lower bound on a flawless success rate tightens to 92.9\%, versus 88.6\% for $n=30$, and the Student-$t$ critical value narrows to $t_{0.975,\,49}=2.010$, producing confidence intervals suited to a Transactions-style submission.

\begin{table}[h]
\centering
\caption{Experimental matrix and measured quantities.}
\label{tab:experiment-matrix}
\begin{tabular}{p{0.21\linewidth} p{0.29\linewidth} p{0.35\linewidth}}
\toprule
\textbf{Experiment} & \textbf{Primary metrics} & \textbf{Data sources} \\
\midrule
Enrollment handshake & enrollment completion latency, enrollment success rate, per-stage goodput & Client-side monotonic timestamps, gateway response path, API visibility of the device \\
Heartbeat supervision & heartbeat observability latency, recent-heartbeat success rate, liveness freshness & Gateway HTTP runtime state, periodic heartbeat timestamps, corroborating Device phase snapshots \\
Application deployment & deploy completion latency, deployment success rate, CRD phase confirmation & Application CRD \texttt{phase} field (authoritative via hardened controller), API server statistics as secondary corroboration \\
Transactional workflow & all-stages success rate, end-to-end latency, end-to-end variability & Per-trial composition of enrollment, heartbeat, and deployment records \\
System snapshot & pod readiness, restart counts, infrastructure metric availability & Kubernetes pod snapshots and platform metrics endpoint \\
\bottomrule
\end{tabular}
\end{table}

The data sources are intentionally multi-layer. For liveness and heartbeat supervision, the gateway HTTP runtime state is authoritative. For deployment, the Application CRD \texttt{phase} field is authoritative because the application controller is active and its reconciliation loop correctly propagates runtime evidence back to the cluster; API statistics are retained as secondary corroboration for reproducibility.

\subsection{Statistical Treatment}
All timing metrics are captured with monotonic client-side clocks to avoid skew from wall-clock adjustments. For continuous variables, the paper reports the mean, sample standard deviation, median, minimum, maximum, 90th/95th/99th percentiles, interquartile range, coefficient of variation, and a 95\% confidence interval for the mean. The confidence interval is computed as
\begin{equation}
\bar{x} \pm t_{0.975,\,n-1}\,\frac{s}{\sqrt{n}},
\end{equation}
where $\bar{x}$ is the sample mean, $s$ is the sample standard deviation, $n$ is the number of valid trials, and $t_{0.975,\,n-1}$ is the Student-$t$ critical value. For success/failure outcomes, the harness reports the success ratio together with a Wilson score interval at 95\% confidence, which is more stable than a naive normal approximation for bounded binary rates.

For throughput-like interpretation, the harness also reports goodput, defined as the number of successful trials divided by the cumulative wall-clock duration of the corresponding measured stage. This value is not intended as a cluster-wide saturation throughput measure; rather, it captures how quickly the evaluated workflow can be completed successfully under the controlled experimental conditions.

The raw collector emits one JSON record per trial and a machine-readable summary for human inspection. This makes it possible to reproduce the reported statistics exactly while also keeping the manuscript readable for reviewers who prefer compact tables and figures over raw traces.

\subsection{Validation Procedure}
Each experimental run follows the same logical procedure.
\begin{enumerate}
    \item Re-establish the runtime prerequisites, including verified API and gateway reachability, TAP/DNAT routing, and active port-forwarding.
    \item Verify that the stable experimental configuration is active, with the gateway, API server, and application controller all running. The controller's reconciliation loop is intentionally active so that the measured path covers the full control plane.
    \item Execute one enrollment attempt and record the elapsed time until the device becomes API-visible and runtime-active.
    \item Query the gateway runtime view of devices and record whether the device remains connected and whether the most recent heartbeat is fresh.
    \item Create and deploy one minimal WASM application through the API server, then inspect runtime deployment statistics for the target device.
    \item Compose the per-trial outcomes into a transactional record, which is successful only if enrollment, heartbeat supervision, and deployment all succeed in the same trial.
\end{enumerate}

This procedure is intentionally traceable: each step leaves evidence either in logs or in Kubernetes resource states, making it possible to reconstruct the run and diagnose deviations.

\subsection{Success Criteria}
A run is considered successful when the following conditions hold simultaneously:
\begin{enumerate}
    \item the TLS handshake and enrollment sequence complete without protocol interruption and the device becomes visible through the northbound API;
    \item the gateway runtime view reports the device as connected and exposes a recent heartbeat timestamp within the acceptance window;
    \item the Application CRD \texttt{phase} field reaches \texttt{Running}, as reconciled by the active application controller; API statistics are retained as secondary corroboration;
    \item the transactional trial is marked successful only when enrollment, heartbeat supervision, and deployment all satisfy their respective criteria in the same trial.
\end{enumerate}

These criteria directly map the system design goals to observable evidence. They validate secure attachment, liveness observability, and remote application activation while avoiding over-reliance on any single eventually consistent status field.

\subsection{Why This Methodology Is Appropriate}
For a systems paper, methodological adequacy depends on whether the collected evidence actually exercises the interaction surfaces claimed as contributions. The present methodology does so because it traverses the full path from cloud-side intent to gateway mediation to edge-side communication and back to northbound observable state. It also captures the distinction between authoritative runtime evidence and noisy control-plane replicas, which is exactly where several of the most important system defects emerged during development.

The companion measurement harness is designed to be run against the same deployment used for validation, so the reported metrics come from the production-like configuration rather than from a synthetic mock. This is important because the dominant risks in RETROSPECT are emergent behaviors at the boundaries between control plane, gateway, and device, and those behaviors only appear when the actual stack is exercised end to end.
